% ---------------------------------------------------------------
% R3.07 SQL dans un langage de programmation -- preambule commun des sujets de TP
% Article clair, fond blanc, accents de couleur coherents avec les slides
% Charte derivee de R5.A.08 Qualite de developpement (meme palette IUT)
% Utilisation :
%   \documentclass[11pt]{article}
%   % ---------------------------------------------------------------
% R3.07 SQL dans un langage de programmation -- preambule commun des sujets de TP
% Article clair, fond blanc, accents de couleur coherents avec les slides
% Charte derivee de R5.A.08 Qualite de developpement (meme palette IUT)
% Utilisation :
%   \documentclass[11pt]{article}
%   % ---------------------------------------------------------------
% R3.07 SQL dans un langage de programmation -- preambule commun des sujets de TP
% Article clair, fond blanc, accents de couleur coherents avec les slides
% Charte derivee de R5.A.08 Qualite de developpement (meme palette IUT)
% Utilisation :
%   \documentclass[11pt]{article}
%   % ---------------------------------------------------------------
% R3.07 SQL dans un langage de programmation -- preambule commun des sujets de TP
% Article clair, fond blanc, accents de couleur coherents avec les slides
% Charte derivee de R5.A.08 Qualite de developpement (meme palette IUT)
% Utilisation :
%   \documentclass[11pt]{article}
%   \input{../../latex/common/r307-tp-preamble.tex}
%   \begin{document}
%   \rtpheader{1}{Titre du TP}
% ---------------------------------------------------------------
\usepackage[T1]{fontenc}
\usepackage[utf8]{inputenc}
\usepackage[french]{babel}
\usepackage[scale=0.92]{tgheros}
\usepackage{tgpagella}
\usepackage[a4paper, margin=2.2cm, headheight=14pt]{geometry}
\usepackage{graphicx}
\usepackage{xcolor}
\usepackage{listings}
\usepackage[skins,listings,breakable]{tcolorbox}
\usepackage{enumitem}
\usepackage{booktabs}
\usepackage{array}
\usepackage{multirow}
\usepackage{tikz}
\usetikzlibrary{shapes, shapes.geometric, arrows, arrows.meta, positioning, calc}
\usepackage[colorlinks=true, linkcolor=qaaccentdark, urlcolor=qaaccentdark]{hyperref}
\usepackage{fancyhdr}
\usepackage{lastpage}
\usepackage{titlesec}
\usepackage{etoolbox}

% ---------------------------------------------------------------
% PALETTE -- version imprimable des couleurs des diapositives
% ---------------------------------------------------------------
\definecolor{qaaccentdark}{HTML}{7A1680}   % violet du logo IUT, fonce pour rester lisible sur blanc
\definecolor{qaaccent2dark}{HTML}{9B4B94}  % violet clair du logo IUT (accent secondaire)
\definecolor{qapanellight}{HTML}{F4F6F8}   % fond de bloc tres clair
\definecolor{qalinelight}{HTML}{D8DEE4}    % filets
\definecolor{qatextdark}{HTML}{1A1D23}     % texte principal
\definecolor{qamuteddark}{HTML}{5B6472}    % texte secondaire
\definecolor{qacodekwl}{HTML}{7A1680}      % code : mots-cles
\definecolor{qacodestrl}{HTML}{2E7D32}     % code : chaines
\definecolor{qacodecoml}{HTML}{6B7280}     % code : commentaires
\definecolor{qacodeannl}{HTML}{B26A00}     % code : annotations / javadoc

\color{qatextdark}

\setlength{\parskip}{4pt}
\setlength{\parindent}{0pt}

% ---------------------------------------------------------------
% TITRES DE SECTION -- texte blanc sur bande de couleur pleine largeur
% (gere aussi les variantes etoilees \section*)
% ---------------------------------------------------------------
\newcommand{\qasectionband}[2]{%
  \noindent\colorbox{#1}{%
    \parbox[c]{\dimexpr\linewidth-2\fboxsep\relax}{%
      \color{white}#2%
    }%
  }%
}
\newcommand{\secformat}[1]{\qasectionband{qaaccentdark}{\bfseries\Large\thesection\quad #1}}
\newcommand{\secformatnn}[1]{\qasectionband{qaaccentdark}{\bfseries\Large #1}}
\newcommand{\subsecformat}[1]{\qasectionband{qaaccent2dark}{\bfseries\large\thesubsection\quad #1}}
\newcommand{\subsecformatnn}[1]{\qasectionband{qaaccent2dark}{\bfseries\large #1}}

\titleformat{\section}{\normalfont}{}{0pt}{\secformat}
\titleformat{name=\section,numberless}{\normalfont}{}{0pt}{\secformatnn}
\titleformat{\subsection}{\normalfont}{}{0pt}{\subsecformat}
\titleformat{name=\subsection,numberless}{\normalfont}{}{0pt}{\subsecformatnn}
\titleformat{\subsubsection}{\normalfont\bfseries\color{qaaccentdark}}{\thesubsubsection}{0.6em}{}
\titlespacing*{\section}{0pt}{14pt}{8pt}
\titlespacing*{\subsection}{0pt}{10pt}{6pt}
\titlespacing*{\subsubsection}{0pt}{8pt}{4pt}

% ---------------------------------------------------------------
% LISTINGS -- theme clair (lstlisting reste utilisable tel quel)
% \lstcode : police + couleur de base, a reutiliser dans les surcharges
% ---------------------------------------------------------------
\newcommand{\lstcode}{\ttfamily\color{qatextdark}}

\lstdefinelanguage{TypeScript}{
  keywords={typeof, new, true, false, catch, function, return, null, switch, var,
    if, in, while, do, else, case, break, const, let, interface, class, extends,
    implements, enum, export, import, this, super, constructor, public, private,
    protected, static, async, await, throw, throws, try, finally, void},
  ndkeywords={boolean, any, number, string, int, long, double},
  sensitive=true,
  comment=[l]{//},
  morecomment=[s]{/*}{*/},
  morestring=[b]',
  morestring=[b]"
}

\lstdefinestyle{r307code}{
  basicstyle=\lstcode\footnotesize,
  keywordstyle=\color{qacodekwl}\bfseries,
  ndkeywordstyle=\color{qaaccent2dark}\bfseries,
  stringstyle=\color{qacodestrl},
  commentstyle=\color{qacodecoml}\itshape,
  morecomment=[s][\color{qacodeannl}]{/**}{*/},
  backgroundcolor=\color{qapanellight},
  rulecolor=\color{qalinelight},
  frame=single,
  framesep=4pt,
  xleftmargin=6pt,
  xrightmargin=6pt,
  numbers=left,
  numberstyle=\tiny\color{qamuteddark},
  numbersep=8pt,
  keepspaces=true,
  columns=fullflexible,
  showstringspaces=false,
  showspaces=false,
  showtabs=false,
  breaklines=true,
  breakatwhitespace=false,
  captionpos=b,
  tabsize=2,
  escapeinside={(*@}{@*)},
}
\lstset{style=r307code}

% ---------------------------------------------------------------
% BOITES : exercice, code, information
% ---------------------------------------------------------------
\tcbset{
  exobox/.style={colback=qapanellight, colframe=qaaccentdark, boxrule=0.6pt,
    arc=2pt, left=8pt, right=8pt, top=6pt, bottom=6pt,
    coltitle=white, colbacktitle=qaaccentdark, fonttitle=\bfseries,
    boxsep=2pt, breakable},
  infobox/.style={colback=white, colframe=qaaccent2dark, boxrule=0.6pt,
    leftrule=3pt, arc=0pt, left=8pt, right=8pt, top=5pt, bottom=5pt,
    coltitle=qaaccent2dark, colbacktitle=white, fonttitle=\bfseries,
    breakable},
  codebox/.style={colback=qatextdark, colframe=qalinelight, boxrule=0.4pt,
    arc=2pt, left=8pt, right=8pt, top=5pt, bottom=5pt, listing only,
    listing options={basicstyle=\ttfamily\small\color{white},
      breaklines=true, showstringspaces=false, columns=fullflexible,
      keepspaces=true, frame=none, numbers=none, backgroundcolor={}}, breakable}
}
\newtcolorbox[auto counter]{exobox}[2][]{exobox,
  title={Exercice~\thetcbcounter\ifstrempty{#2}{}{ --- #2}}, #1}
\newtcolorbox{infobox}[1][\`A retenir]{infobox, title={#1}}
\newtcblisting{codebox}{codebox}

% ---------------------------------------------------------------
% EN-TETE DE PAGE (fil d'Ariane + pagination)
% ---------------------------------------------------------------
\newcommand{\qaheadtitlevalue}{}
\newcommand{\qaheadtitle}[1]{\renewcommand{\qaheadtitlevalue}{#1}}

\pagestyle{fancy}
\fancyhf{}
\renewcommand{\headrulewidth}{0.4pt}
\fancyhead[L]{\small\color{qamuteddark}R3.07 --- \qaheadtitlevalue}
\fancyhead[R]{\small\color{qamuteddark}\thepage\,/\,\pageref{LastPage}}
\fancyfoot[C]{\footnotesize\color{qamuteddark}IUT de Vannes --- Universit\'e Bretagne Sud}
\fancypagestyle{plain}{\fancyhf{}\renewcommand{\headrulewidth}{0pt}%
  \fancyfoot[C]{\footnotesize\color{qamuteddark}IUT de Vannes --- Universit\'e Bretagne Sud}}

% ---------------------------------------------------------------
% BANDEAU DE TITRE (pas une page dediee, pour rester compact)
% \rtpheader{numero ou intitule court}{titre du TP}
% ---------------------------------------------------------------
\newcommand{\ryear}{2026/2027}
\newcommand{\rtpheader}[2]{%
  \thispagestyle{plain}%
  \begin{minipage}[t]{0.68\textwidth}
    {\small\color{qamuteddark}
      Universit\'e Bretagne Sud\\
      IUT de Vannes\\
      D\'epartement Informatique --- \ryear}
  \end{minipage}%
  \hfill
  \begin{minipage}[t]{0.28\textwidth}
    \raggedleft
    \includegraphics[width=2.6cm]{iut_logo.jpg}
  \end{minipage}

  \vspace{6mm}
  \begin{center}
    {\large\bfseries\color{qatextdark} BUT 2 --- R3.07. SQL dans un langage de programmation}\\[2mm]
    {\normalsize\color{qamuteddark} TD/TP #1}\\[1mm]
    {\Large\bfseries\color{qaaccentdark} #2}
  \end{center}
  \vspace{4mm}
  \qaheadtitle{#2}
}

%   \begin{document}
%   \rtpheader{1}{Titre du TP}
% ---------------------------------------------------------------
\usepackage[T1]{fontenc}
\usepackage[utf8]{inputenc}
\usepackage[french]{babel}
\usepackage[scale=0.92]{tgheros}
\usepackage{tgpagella}
\usepackage[a4paper, margin=2.2cm, headheight=14pt]{geometry}
\usepackage{graphicx}
\usepackage{xcolor}
\usepackage{listings}
\usepackage[skins,listings,breakable]{tcolorbox}
\usepackage{enumitem}
\usepackage{booktabs}
\usepackage{array}
\usepackage{multirow}
\usepackage{tikz}
\usetikzlibrary{shapes, shapes.geometric, arrows, arrows.meta, positioning, calc}
\usepackage[colorlinks=true, linkcolor=qaaccentdark, urlcolor=qaaccentdark]{hyperref}
\usepackage{fancyhdr}
\usepackage{lastpage}
\usepackage{titlesec}
\usepackage{etoolbox}

% ---------------------------------------------------------------
% PALETTE -- version imprimable des couleurs des diapositives
% ---------------------------------------------------------------
\definecolor{qaaccentdark}{HTML}{7A1680}   % violet du logo IUT, fonce pour rester lisible sur blanc
\definecolor{qaaccent2dark}{HTML}{9B4B94}  % violet clair du logo IUT (accent secondaire)
\definecolor{qapanellight}{HTML}{F4F6F8}   % fond de bloc tres clair
\definecolor{qalinelight}{HTML}{D8DEE4}    % filets
\definecolor{qatextdark}{HTML}{1A1D23}     % texte principal
\definecolor{qamuteddark}{HTML}{5B6472}    % texte secondaire
\definecolor{qacodekwl}{HTML}{7A1680}      % code : mots-cles
\definecolor{qacodestrl}{HTML}{2E7D32}     % code : chaines
\definecolor{qacodecoml}{HTML}{6B7280}     % code : commentaires
\definecolor{qacodeannl}{HTML}{B26A00}     % code : annotations / javadoc

\color{qatextdark}

\setlength{\parskip}{4pt}
\setlength{\parindent}{0pt}

% ---------------------------------------------------------------
% TITRES DE SECTION -- texte blanc sur bande de couleur pleine largeur
% (gere aussi les variantes etoilees \section*)
% ---------------------------------------------------------------
\newcommand{\qasectionband}[2]{%
  \noindent\colorbox{#1}{%
    \parbox[c]{\dimexpr\linewidth-2\fboxsep\relax}{%
      \color{white}#2%
    }%
  }%
}
\newcommand{\secformat}[1]{\qasectionband{qaaccentdark}{\bfseries\Large\thesection\quad #1}}
\newcommand{\secformatnn}[1]{\qasectionband{qaaccentdark}{\bfseries\Large #1}}
\newcommand{\subsecformat}[1]{\qasectionband{qaaccent2dark}{\bfseries\large\thesubsection\quad #1}}
\newcommand{\subsecformatnn}[1]{\qasectionband{qaaccent2dark}{\bfseries\large #1}}

\titleformat{\section}{\normalfont}{}{0pt}{\secformat}
\titleformat{name=\section,numberless}{\normalfont}{}{0pt}{\secformatnn}
\titleformat{\subsection}{\normalfont}{}{0pt}{\subsecformat}
\titleformat{name=\subsection,numberless}{\normalfont}{}{0pt}{\subsecformatnn}
\titleformat{\subsubsection}{\normalfont\bfseries\color{qaaccentdark}}{\thesubsubsection}{0.6em}{}
\titlespacing*{\section}{0pt}{14pt}{8pt}
\titlespacing*{\subsection}{0pt}{10pt}{6pt}
\titlespacing*{\subsubsection}{0pt}{8pt}{4pt}

% ---------------------------------------------------------------
% LISTINGS -- theme clair (lstlisting reste utilisable tel quel)
% \lstcode : police + couleur de base, a reutiliser dans les surcharges
% ---------------------------------------------------------------
\newcommand{\lstcode}{\ttfamily\color{qatextdark}}

\lstdefinelanguage{TypeScript}{
  keywords={typeof, new, true, false, catch, function, return, null, switch, var,
    if, in, while, do, else, case, break, const, let, interface, class, extends,
    implements, enum, export, import, this, super, constructor, public, private,
    protected, static, async, await, throw, throws, try, finally, void},
  ndkeywords={boolean, any, number, string, int, long, double},
  sensitive=true,
  comment=[l]{//},
  morecomment=[s]{/*}{*/},
  morestring=[b]',
  morestring=[b]"
}

\lstdefinestyle{r307code}{
  basicstyle=\lstcode\footnotesize,
  keywordstyle=\color{qacodekwl}\bfseries,
  ndkeywordstyle=\color{qaaccent2dark}\bfseries,
  stringstyle=\color{qacodestrl},
  commentstyle=\color{qacodecoml}\itshape,
  morecomment=[s][\color{qacodeannl}]{/**}{*/},
  backgroundcolor=\color{qapanellight},
  rulecolor=\color{qalinelight},
  frame=single,
  framesep=4pt,
  xleftmargin=6pt,
  xrightmargin=6pt,
  numbers=left,
  numberstyle=\tiny\color{qamuteddark},
  numbersep=8pt,
  keepspaces=true,
  columns=fullflexible,
  showstringspaces=false,
  showspaces=false,
  showtabs=false,
  breaklines=true,
  breakatwhitespace=false,
  captionpos=b,
  tabsize=2,
  escapeinside={(*@}{@*)},
}
\lstset{style=r307code}

% ---------------------------------------------------------------
% BOITES : exercice, code, information
% ---------------------------------------------------------------
\tcbset{
  exobox/.style={colback=qapanellight, colframe=qaaccentdark, boxrule=0.6pt,
    arc=2pt, left=8pt, right=8pt, top=6pt, bottom=6pt,
    coltitle=white, colbacktitle=qaaccentdark, fonttitle=\bfseries,
    boxsep=2pt, breakable},
  infobox/.style={colback=white, colframe=qaaccent2dark, boxrule=0.6pt,
    leftrule=3pt, arc=0pt, left=8pt, right=8pt, top=5pt, bottom=5pt,
    coltitle=qaaccent2dark, colbacktitle=white, fonttitle=\bfseries,
    breakable},
  codebox/.style={colback=qatextdark, colframe=qalinelight, boxrule=0.4pt,
    arc=2pt, left=8pt, right=8pt, top=5pt, bottom=5pt, listing only,
    listing options={basicstyle=\ttfamily\small\color{white},
      breaklines=true, showstringspaces=false, columns=fullflexible,
      keepspaces=true, frame=none, numbers=none, backgroundcolor={}}, breakable}
}
\newtcolorbox[auto counter]{exobox}[2][]{exobox,
  title={Exercice~\thetcbcounter\ifstrempty{#2}{}{ --- #2}}, #1}
\newtcolorbox{infobox}[1][\`A retenir]{infobox, title={#1}}
\newtcblisting{codebox}{codebox}

% ---------------------------------------------------------------
% EN-TETE DE PAGE (fil d'Ariane + pagination)
% ---------------------------------------------------------------
\newcommand{\qaheadtitlevalue}{}
\newcommand{\qaheadtitle}[1]{\renewcommand{\qaheadtitlevalue}{#1}}

\pagestyle{fancy}
\fancyhf{}
\renewcommand{\headrulewidth}{0.4pt}
\fancyhead[L]{\small\color{qamuteddark}R3.07 --- \qaheadtitlevalue}
\fancyhead[R]{\small\color{qamuteddark}\thepage\,/\,\pageref{LastPage}}
\fancyfoot[C]{\footnotesize\color{qamuteddark}IUT de Vannes --- Universit\'e Bretagne Sud}
\fancypagestyle{plain}{\fancyhf{}\renewcommand{\headrulewidth}{0pt}%
  \fancyfoot[C]{\footnotesize\color{qamuteddark}IUT de Vannes --- Universit\'e Bretagne Sud}}

% ---------------------------------------------------------------
% BANDEAU DE TITRE (pas une page dediee, pour rester compact)
% \rtpheader{numero ou intitule court}{titre du TP}
% ---------------------------------------------------------------
\newcommand{\ryear}{2026/2027}
\newcommand{\rtpheader}[2]{%
  \thispagestyle{plain}%
  \begin{minipage}[t]{0.68\textwidth}
    {\small\color{qamuteddark}
      Universit\'e Bretagne Sud\\
      IUT de Vannes\\
      D\'epartement Informatique --- \ryear}
  \end{minipage}%
  \hfill
  \begin{minipage}[t]{0.28\textwidth}
    \raggedleft
    \includegraphics[width=2.6cm]{iut_logo.jpg}
  \end{minipage}

  \vspace{6mm}
  \begin{center}
    {\large\bfseries\color{qatextdark} BUT 2 --- R3.07. SQL dans un langage de programmation}\\[2mm]
    {\normalsize\color{qamuteddark} TD/TP #1}\\[1mm]
    {\Large\bfseries\color{qaaccentdark} #2}
  \end{center}
  \vspace{4mm}
  \qaheadtitle{#2}
}

%   \begin{document}
%   \rtpheader{1}{Titre du TP}
% ---------------------------------------------------------------
\usepackage[T1]{fontenc}
\usepackage[utf8]{inputenc}
\usepackage[french]{babel}
\usepackage[scale=0.92]{tgheros}
\usepackage{tgpagella}
\usepackage[a4paper, margin=2.2cm, headheight=14pt]{geometry}
\usepackage{graphicx}
\usepackage{xcolor}
\usepackage{listings}
\usepackage[skins,listings,breakable]{tcolorbox}
\usepackage{enumitem}
\usepackage{booktabs}
\usepackage{array}
\usepackage{multirow}
\usepackage{tikz}
\usetikzlibrary{shapes, shapes.geometric, arrows, arrows.meta, positioning, calc}
\usepackage[colorlinks=true, linkcolor=qaaccentdark, urlcolor=qaaccentdark]{hyperref}
\usepackage{fancyhdr}
\usepackage{lastpage}
\usepackage{titlesec}
\usepackage{etoolbox}

% ---------------------------------------------------------------
% PALETTE -- version imprimable des couleurs des diapositives
% ---------------------------------------------------------------
\definecolor{qaaccentdark}{HTML}{7A1680}   % violet du logo IUT, fonce pour rester lisible sur blanc
\definecolor{qaaccent2dark}{HTML}{9B4B94}  % violet clair du logo IUT (accent secondaire)
\definecolor{qapanellight}{HTML}{F4F6F8}   % fond de bloc tres clair
\definecolor{qalinelight}{HTML}{D8DEE4}    % filets
\definecolor{qatextdark}{HTML}{1A1D23}     % texte principal
\definecolor{qamuteddark}{HTML}{5B6472}    % texte secondaire
\definecolor{qacodekwl}{HTML}{7A1680}      % code : mots-cles
\definecolor{qacodestrl}{HTML}{2E7D32}     % code : chaines
\definecolor{qacodecoml}{HTML}{6B7280}     % code : commentaires
\definecolor{qacodeannl}{HTML}{B26A00}     % code : annotations / javadoc

\color{qatextdark}

\setlength{\parskip}{4pt}
\setlength{\parindent}{0pt}

% ---------------------------------------------------------------
% TITRES DE SECTION -- texte blanc sur bande de couleur pleine largeur
% (gere aussi les variantes etoilees \section*)
% ---------------------------------------------------------------
\newcommand{\qasectionband}[2]{%
  \noindent\colorbox{#1}{%
    \parbox[c]{\dimexpr\linewidth-2\fboxsep\relax}{%
      \color{white}#2%
    }%
  }%
}
\newcommand{\secformat}[1]{\qasectionband{qaaccentdark}{\bfseries\Large\thesection\quad #1}}
\newcommand{\secformatnn}[1]{\qasectionband{qaaccentdark}{\bfseries\Large #1}}
\newcommand{\subsecformat}[1]{\qasectionband{qaaccent2dark}{\bfseries\large\thesubsection\quad #1}}
\newcommand{\subsecformatnn}[1]{\qasectionband{qaaccent2dark}{\bfseries\large #1}}

\titleformat{\section}{\normalfont}{}{0pt}{\secformat}
\titleformat{name=\section,numberless}{\normalfont}{}{0pt}{\secformatnn}
\titleformat{\subsection}{\normalfont}{}{0pt}{\subsecformat}
\titleformat{name=\subsection,numberless}{\normalfont}{}{0pt}{\subsecformatnn}
\titleformat{\subsubsection}{\normalfont\bfseries\color{qaaccentdark}}{\thesubsubsection}{0.6em}{}
\titlespacing*{\section}{0pt}{14pt}{8pt}
\titlespacing*{\subsection}{0pt}{10pt}{6pt}
\titlespacing*{\subsubsection}{0pt}{8pt}{4pt}

% ---------------------------------------------------------------
% LISTINGS -- theme clair (lstlisting reste utilisable tel quel)
% \lstcode : police + couleur de base, a reutiliser dans les surcharges
% ---------------------------------------------------------------
\newcommand{\lstcode}{\ttfamily\color{qatextdark}}

\lstdefinelanguage{TypeScript}{
  keywords={typeof, new, true, false, catch, function, return, null, switch, var,
    if, in, while, do, else, case, break, const, let, interface, class, extends,
    implements, enum, export, import, this, super, constructor, public, private,
    protected, static, async, await, throw, throws, try, finally, void},
  ndkeywords={boolean, any, number, string, int, long, double},
  sensitive=true,
  comment=[l]{//},
  morecomment=[s]{/*}{*/},
  morestring=[b]',
  morestring=[b]"
}

\lstdefinestyle{r307code}{
  basicstyle=\lstcode\footnotesize,
  keywordstyle=\color{qacodekwl}\bfseries,
  ndkeywordstyle=\color{qaaccent2dark}\bfseries,
  stringstyle=\color{qacodestrl},
  commentstyle=\color{qacodecoml}\itshape,
  morecomment=[s][\color{qacodeannl}]{/**}{*/},
  backgroundcolor=\color{qapanellight},
  rulecolor=\color{qalinelight},
  frame=single,
  framesep=4pt,
  xleftmargin=6pt,
  xrightmargin=6pt,
  numbers=left,
  numberstyle=\tiny\color{qamuteddark},
  numbersep=8pt,
  keepspaces=true,
  columns=fullflexible,
  showstringspaces=false,
  showspaces=false,
  showtabs=false,
  breaklines=true,
  breakatwhitespace=false,
  captionpos=b,
  tabsize=2,
  escapeinside={(*@}{@*)},
}
\lstset{style=r307code}

% ---------------------------------------------------------------
% BOITES : exercice, code, information
% ---------------------------------------------------------------
\tcbset{
  exobox/.style={colback=qapanellight, colframe=qaaccentdark, boxrule=0.6pt,
    arc=2pt, left=8pt, right=8pt, top=6pt, bottom=6pt,
    coltitle=white, colbacktitle=qaaccentdark, fonttitle=\bfseries,
    boxsep=2pt, breakable},
  infobox/.style={colback=white, colframe=qaaccent2dark, boxrule=0.6pt,
    leftrule=3pt, arc=0pt, left=8pt, right=8pt, top=5pt, bottom=5pt,
    coltitle=qaaccent2dark, colbacktitle=white, fonttitle=\bfseries,
    breakable},
  codebox/.style={colback=qatextdark, colframe=qalinelight, boxrule=0.4pt,
    arc=2pt, left=8pt, right=8pt, top=5pt, bottom=5pt, listing only,
    listing options={basicstyle=\ttfamily\small\color{white},
      breaklines=true, showstringspaces=false, columns=fullflexible,
      keepspaces=true, frame=none, numbers=none, backgroundcolor={}}, breakable}
}
\newtcolorbox[auto counter]{exobox}[2][]{exobox,
  title={Exercice~\thetcbcounter\ifstrempty{#2}{}{ --- #2}}, #1}
\newtcolorbox{infobox}[1][\`A retenir]{infobox, title={#1}}
\newtcblisting{codebox}{codebox}

% ---------------------------------------------------------------
% EN-TETE DE PAGE (fil d'Ariane + pagination)
% ---------------------------------------------------------------
\newcommand{\qaheadtitlevalue}{}
\newcommand{\qaheadtitle}[1]{\renewcommand{\qaheadtitlevalue}{#1}}

\pagestyle{fancy}
\fancyhf{}
\renewcommand{\headrulewidth}{0.4pt}
\fancyhead[L]{\small\color{qamuteddark}R3.07 --- \qaheadtitlevalue}
\fancyhead[R]{\small\color{qamuteddark}\thepage\,/\,\pageref{LastPage}}
\fancyfoot[C]{\footnotesize\color{qamuteddark}IUT de Vannes --- Universit\'e Bretagne Sud}
\fancypagestyle{plain}{\fancyhf{}\renewcommand{\headrulewidth}{0pt}%
  \fancyfoot[C]{\footnotesize\color{qamuteddark}IUT de Vannes --- Universit\'e Bretagne Sud}}

% ---------------------------------------------------------------
% BANDEAU DE TITRE (pas une page dediee, pour rester compact)
% \rtpheader{numero ou intitule court}{titre du TP}
% ---------------------------------------------------------------
\newcommand{\ryear}{2026/2027}
\newcommand{\rtpheader}[2]{%
  \thispagestyle{plain}%
  \begin{minipage}[t]{0.68\textwidth}
    {\small\color{qamuteddark}
      Universit\'e Bretagne Sud\\
      IUT de Vannes\\
      D\'epartement Informatique --- \ryear}
  \end{minipage}%
  \hfill
  \begin{minipage}[t]{0.28\textwidth}
    \raggedleft
    \includegraphics[width=2.6cm]{iut_logo.jpg}
  \end{minipage}

  \vspace{6mm}
  \begin{center}
    {\large\bfseries\color{qatextdark} BUT 2 --- R3.07. SQL dans un langage de programmation}\\[2mm]
    {\normalsize\color{qamuteddark} TD/TP #1}\\[1mm]
    {\Large\bfseries\color{qaaccentdark} #2}
  \end{center}
  \vspace{4mm}
  \qaheadtitle{#2}
}

%   \begin{document}
%   \rtpheader{1}{Titre du TP}
% ---------------------------------------------------------------
\usepackage[T1]{fontenc}
\usepackage[utf8]{inputenc}
\usepackage[french]{babel}
\usepackage[scale=0.92]{tgheros}
\usepackage{tgpagella}
\usepackage[a4paper, margin=2.2cm, headheight=14pt]{geometry}
\usepackage{graphicx}
\usepackage{xcolor}
\usepackage{listings}
\usepackage[skins,listings,breakable]{tcolorbox}
\usepackage{enumitem}
\usepackage{booktabs}
\usepackage{array}
\usepackage{multirow}
\usepackage{tikz}
\usetikzlibrary{shapes, shapes.geometric, arrows, arrows.meta, positioning, calc}
\usepackage[colorlinks=true, linkcolor=qaaccentdark, urlcolor=qaaccentdark]{hyperref}
\usepackage{fancyhdr}
\usepackage{lastpage}
\usepackage{titlesec}
\usepackage{etoolbox}

% ---------------------------------------------------------------
% PALETTE -- version imprimable des couleurs des diapositives
% ---------------------------------------------------------------
\definecolor{qaaccentdark}{HTML}{7A1680}   % violet du logo IUT, fonce pour rester lisible sur blanc
\definecolor{qaaccent2dark}{HTML}{9B4B94}  % violet clair du logo IUT (accent secondaire)
\definecolor{qapanellight}{HTML}{F4F6F8}   % fond de bloc tres clair
\definecolor{qalinelight}{HTML}{D8DEE4}    % filets
\definecolor{qatextdark}{HTML}{1A1D23}     % texte principal
\definecolor{qamuteddark}{HTML}{5B6472}    % texte secondaire
\definecolor{qacodekwl}{HTML}{7A1680}      % code : mots-cles
\definecolor{qacodestrl}{HTML}{2E7D32}     % code : chaines
\definecolor{qacodecoml}{HTML}{6B7280}     % code : commentaires
\definecolor{qacodeannl}{HTML}{B26A00}     % code : annotations / javadoc

\color{qatextdark}

\setlength{\parskip}{4pt}
\setlength{\parindent}{0pt}

% ---------------------------------------------------------------
% TITRES DE SECTION -- texte blanc sur bande de couleur pleine largeur
% (gere aussi les variantes etoilees \section*)
% ---------------------------------------------------------------
\newcommand{\qasectionband}[2]{%
  \noindent\colorbox{#1}{%
    \parbox[c]{\dimexpr\linewidth-2\fboxsep\relax}{%
      \color{white}#2%
    }%
  }%
}
\newcommand{\secformat}[1]{\qasectionband{qaaccentdark}{\bfseries\Large\thesection\quad #1}}
\newcommand{\secformatnn}[1]{\qasectionband{qaaccentdark}{\bfseries\Large #1}}
\newcommand{\subsecformat}[1]{\qasectionband{qaaccent2dark}{\bfseries\large\thesubsection\quad #1}}
\newcommand{\subsecformatnn}[1]{\qasectionband{qaaccent2dark}{\bfseries\large #1}}

\titleformat{\section}{\normalfont}{}{0pt}{\secformat}
\titleformat{name=\section,numberless}{\normalfont}{}{0pt}{\secformatnn}
\titleformat{\subsection}{\normalfont}{}{0pt}{\subsecformat}
\titleformat{name=\subsection,numberless}{\normalfont}{}{0pt}{\subsecformatnn}
\titleformat{\subsubsection}{\normalfont\bfseries\color{qaaccentdark}}{\thesubsubsection}{0.6em}{}
\titlespacing*{\section}{0pt}{14pt}{8pt}
\titlespacing*{\subsection}{0pt}{10pt}{6pt}
\titlespacing*{\subsubsection}{0pt}{8pt}{4pt}

% ---------------------------------------------------------------
% LISTINGS -- theme clair (lstlisting reste utilisable tel quel)
% \lstcode : police + couleur de base, a reutiliser dans les surcharges
% ---------------------------------------------------------------
\newcommand{\lstcode}{\ttfamily\color{qatextdark}}

\lstdefinelanguage{TypeScript}{
  keywords={typeof, new, true, false, catch, function, return, null, switch, var,
    if, in, while, do, else, case, break, const, let, interface, class, extends,
    implements, enum, export, import, this, super, constructor, public, private,
    protected, static, async, await, throw, throws, try, finally, void},
  ndkeywords={boolean, any, number, string, int, long, double},
  sensitive=true,
  comment=[l]{//},
  morecomment=[s]{/*}{*/},
  morestring=[b]',
  morestring=[b]"
}

\lstdefinestyle{r307code}{
  basicstyle=\lstcode\footnotesize,
  keywordstyle=\color{qacodekwl}\bfseries,
  ndkeywordstyle=\color{qaaccent2dark}\bfseries,
  stringstyle=\color{qacodestrl},
  commentstyle=\color{qacodecoml}\itshape,
  morecomment=[s][\color{qacodeannl}]{/**}{*/},
  backgroundcolor=\color{qapanellight},
  rulecolor=\color{qalinelight},
  frame=single,
  framesep=4pt,
  xleftmargin=6pt,
  xrightmargin=6pt,
  numbers=left,
  numberstyle=\tiny\color{qamuteddark},
  numbersep=8pt,
  keepspaces=true,
  columns=fullflexible,
  showstringspaces=false,
  showspaces=false,
  showtabs=false,
  breaklines=true,
  breakatwhitespace=false,
  captionpos=b,
  tabsize=2,
  escapeinside={(*@}{@*)},
}
\lstset{style=r307code}

% ---------------------------------------------------------------
% BOITES : exercice, code, information
% ---------------------------------------------------------------
\tcbset{
  exobox/.style={colback=qapanellight, colframe=qaaccentdark, boxrule=0.6pt,
    arc=2pt, left=8pt, right=8pt, top=6pt, bottom=6pt,
    coltitle=white, colbacktitle=qaaccentdark, fonttitle=\bfseries,
    boxsep=2pt, breakable},
  infobox/.style={colback=white, colframe=qaaccent2dark, boxrule=0.6pt,
    leftrule=3pt, arc=0pt, left=8pt, right=8pt, top=5pt, bottom=5pt,
    coltitle=qaaccent2dark, colbacktitle=white, fonttitle=\bfseries,
    breakable},
  codebox/.style={colback=qatextdark, colframe=qalinelight, boxrule=0.4pt,
    arc=2pt, left=8pt, right=8pt, top=5pt, bottom=5pt, listing only,
    listing options={basicstyle=\ttfamily\small\color{white},
      breaklines=true, showstringspaces=false, columns=fullflexible,
      keepspaces=true, frame=none, numbers=none, backgroundcolor={}}, breakable}
}
\newtcolorbox[auto counter]{exobox}[2][]{exobox,
  title={Exercice~\thetcbcounter\ifstrempty{#2}{}{ --- #2}}, #1}
\newtcolorbox{infobox}[1][\`A retenir]{infobox, title={#1}}
\newtcblisting{codebox}{codebox}

% ---------------------------------------------------------------
% EN-TETE DE PAGE (fil d'Ariane + pagination)
% ---------------------------------------------------------------
\newcommand{\qaheadtitlevalue}{}
\newcommand{\qaheadtitle}[1]{\renewcommand{\qaheadtitlevalue}{#1}}

\pagestyle{fancy}
\fancyhf{}
\renewcommand{\headrulewidth}{0.4pt}
\fancyhead[L]{\small\color{qamuteddark}R3.07 --- \qaheadtitlevalue}
\fancyhead[R]{\small\color{qamuteddark}\thepage\,/\,\pageref{LastPage}}
\fancyfoot[C]{\footnotesize\color{qamuteddark}IUT de Vannes --- Universit\'e Bretagne Sud}
\fancypagestyle{plain}{\fancyhf{}\renewcommand{\headrulewidth}{0pt}%
  \fancyfoot[C]{\footnotesize\color{qamuteddark}IUT de Vannes --- Universit\'e Bretagne Sud}}

% ---------------------------------------------------------------
% BANDEAU DE TITRE (pas une page dediee, pour rester compact)
% \rtpheader{numero ou intitule court}{titre du TP}
% ---------------------------------------------------------------
\newcommand{\ryear}{2026/2027}
\newcommand{\rtpheader}[2]{%
  \thispagestyle{plain}%
  \begin{minipage}[t]{0.68\textwidth}
    {\small\color{qamuteddark}
      Universit\'e Bretagne Sud\\
      IUT de Vannes\\
      D\'epartement Informatique --- \ryear}
  \end{minipage}%
  \hfill
  \begin{minipage}[t]{0.28\textwidth}
    \raggedleft
    \includegraphics[width=2.6cm]{iut_logo.jpg}
  \end{minipage}

  \vspace{6mm}
  \begin{center}
    {\large\bfseries\color{qatextdark} BUT 2 --- R3.07. SQL dans un langage de programmation}\\[2mm]
    {\normalsize\color{qamuteddark} TD/TP #1}\\[1mm]
    {\Large\bfseries\color{qaaccentdark} #2}
  \end{center}
  \vspace{4mm}
  \qaheadtitle{#2}
}
